\section{Recommendation System}\label{sec:recommendation}

\subsection{Related work}
\label{subsec:rs_related_work}

A large language model (LLM) is inherently designed to predict the next token in a sequence, which makes it well-suited for sequential recommendation tasks when trained on historical patterns of user behavior \cite{P52022, li2023text}. Rather than retraining an LLM from scratch, an alternative approach leverages pre-trained LLMs for recommendation. For instance, zero-shot learning can be employed when the model utilizes textual descriptions (e.g., song titles or movie names) rather than unique object identifiers, enabling it to draw upon its prior knowledge \cite{sanner2023large}. Additionally, by including examples of user-item interactions within the LLM's prompt, it is possible to elicit meaningful recommendations through in-context learning \cite{xu2024prompting} or via some lightweight instruction tuning \cite{bao2023tallrec}.

There exist also specific ways of encoding the input in an LLM, for example, how to encode graph structures \cite{tang2024graphgpt}. The interested reader can also find a summary of recent advances of LLM-based approaches for recommender systems in \cite{deldjoo2024review, zhao2024recommender}.

In our work, we focus on an approach that is simple to implement and can be combined with powerful existing recommender architectures, namely using the LLM as an embedding to represent the items of a recommendation system. For instance, when the item is a book, we can concatenate its text attributes (title, author, and abstract) and use the LLM to map this information into a new representation space.
This representation will be used in two ways: a) as a representation of the items, which can later be used for a KNN (K-Nearest-Neighbor) approach for finding appropriate recommendations, b) serve as the initialization for embeddings in a graph neural network, specifically the LightGCN architecture \cite{lightgcn}, which is known for its speed and accuracy in predictions. This hybrid approach allows us to combine the advanced text representation capabilities of LLMs with the efficacy of graph neural networks, while maintaining modularity. Additionally, it avoids the need for costly retraining of the LLM, making the methodology easy to code and deploy, and also resource-efficient.


\textcolor{red}{Embedding techniques in graph-based recommendation methods have been studied extensively to integrate content-based representations into collaborative filtering architectures \cite{koren2021advances,lightgcn,bpr}. Pairwise learning-to-rank formulations such as Bayesian Personalized Ranking (BPR) \cite{bpr} and matrix factorization methods for implicit feedback (ALS, LMF) \cite{als1,als2,lmf} optimize item orderings from historical interactions, but suffer under severe cold-start sparsity when items lack interaction histories. Recent work shows that initializing graph neural networks with dense representations extracted from pre-trained LLMs bridges content-based semantics and high-order collaborative signals \cite{deldjoo2024review,zhao2024recommender}.}

\subsection{Our approach}
\label{subsec:rs_our_approach}

Building on our models for difficulty estimation and simplification, we introduce a recommender system that uses these outputs to suggest suitable content for language learners. In language learning applications, a recommender system is essential for tailoring content to match users’ proficiency levels and personal interests.

When the system identifies (or when the user explicitly declares) their proficiency level, it filters and ranks content that matches both the user’s level and topical interests (e.g., sports or technology), based on data from the user’s historical preferences.

Various approaches exist for recommender systems, including collaborative filtering, content-based filtering, knowledge-based recommendation, and hybrid combinations. Collaborative filtering algorithms, in particular, assume that accurate recommendations can be made by using the collective preferences of similar, like-minded users \cite{koren2021advances}. \textcolor{red}{Our methodological contribution in Section~\ref{sec:recommendation} consists of injecting dense semantic representations from pre-trained LLMs directly into both neighborhood-based (KNN) and graph-convolutional (LightGCN) collaborative filtering architectures without requiring end-to-end LLM fine-tuning.}

\subsubsection*{\textcolor{red}{1. LLM Item and User Embedding Generation (\autoref{fig:init_embeddings})}}
\textcolor{red}{Let $\mathcal{U}$ denote the set of users and $\mathcal{I}$ the set of candidate reading items. As depicted in \autoref{fig:init_embeddings}, each item $i \in \mathcal{I}$ is associated with a textual document or concatenated metadata string $d_i$ (such as the full article text in \texttt{Zeeguu}, or the title, summary, and metadata in \texttt{Goodreads}, \texttt{ml-100k}, and \texttt{Tomplay}). We pass $d_i$ through a pre-trained language model encoder $\phi_{\text{LLM}}$ and apply mean pooling over its contextualized token hidden states $\mathbf{h}_{i,1}, \dots, \mathbf{h}_{i,|T_i|}$ to obtain a $d$-dimensional item embedding $\mathbf{e}_i^{(0)} \in \mathbb{R}^d$:}
{\color{red}\begin{equation}
    \mathbf{e}_i^{(0)} = \phi_{\text{LLM}}(d_i) = \frac{1}{|T_i|} \sum_{t=1}^{|T_i|} \mathbf{h}_{i,t} \in \mathbb{R}^d,
    \label{eq:item_embedding_init}
\end{equation}}
\textcolor{red}{where $d = 768$ for our multilingual RoBERTa sentence encoder (\texttt{paraphrase-xlm-r-multilingual-v1}, $|T_i| \le 128$) and $d = 1{,}536$ for OpenAI \texttt{text-embedding-ada-002} ($|T_i| \le 8{,}000$). To represent each user $u \in \mathcal{U}$ in the same semantic space $\mathbb{R}^d$ without requiring explicit demographic profiles, we aggregate the embeddings of the $n$ most recent items $\mathcal{H}_u^{(n)} = \{i_{u,1}, \dots, i_{u,n}\}$ with which user $u$ has positively interacted:}
{\color{red}\begin{equation}
    \mathbf{e}_u^{(0)} = \frac{1}{|\mathcal{H}_u^{(n)}|} \sum_{i \in \mathcal{H}_u^{(n)}} \mathbf{e}_i^{(0)} \in \mathbb{R}^d.
    \label{eq:user_embedding_init}
\end{equation}}

\begin{figure}[ht!]
    \centering
    \includegraphics[width=\linewidth]{figures/init_embeddings.png}
    \caption{Generating item embeddings with large language models (LLMs) and combining them to construct user embeddings.}
    \label{fig:init_embeddings}
\end{figure}

\begin{itemize}
    \item \textbf{KNN-based \textcolor{red}{(\autoref{fig:recommend_with_knn})}}: \textcolor{red}{In the user-based K-Nearest Neighbors (KNN) method illustrated in \autoref{fig:recommend_with_knn}, we measure the semantic preference similarity between target user $u$ and every other user $v \in \mathcal{U} \setminus \{u\}$ via cosine similarity in the LLM embedding space:}
    {\color{red}\begin{equation}
        \text{sim}(u, v) = \frac{\mathbf{e}_u^{(0)} \cdot \mathbf{e}_v^{(0)}}{\|\mathbf{e}_u^{(0)}\|_2 \, \|\mathbf{e}_v^{(0)}\|_2}.
        \label{eq:knn_user_sim}
    \end{equation}}
    \textcolor{red}{Let $\mathcal{N}_k(u) \subset \mathcal{U} \setminus \{u\}$ denote the set of $k$ nearest neighbor users maximizing $\text{sim}(u, v)$. For any candidate item $i \notin \mathcal{H}_u$ whose estimated (or simplified) CEFR difficulty matches the learner's proficiency level, the predicted preference score $\hat{y}_{u,i}^{\text{KNN}}$ is computed as the similarity-weighted sum of neighbor interactions $r_{v,i}$:}
    {\color{red}\begin{equation}
        \hat{y}_{u,i}^{\text{KNN}} = \sum_{v \in \mathcal{N}_k(u)} \text{sim}(u, v) \cdot r_{v,i}.
        \label{eq:knn_score}
    \end{equation}}

    \begin{figure}[ht!]
        \centering
        \includegraphics[width=0.5\linewidth]{figures/recommend_with_knn.png}
        \caption{Selecting recommended items for a user based on similar users’ preferences using the KNN approach.}
        \label{fig:recommend_with_knn}
    \end{figure}
 
    \item \textbf{Graph-based \textcolor{red}{(\autoref{fig:recommend_with_lightgcn})}}: We adapt the LightGCN architecture \cite{lightgcn} \textcolor{red}{on the bipartite user-item interaction graph $\mathcal{G} = (\mathcal{U} \cup \mathcal{I}, \mathcal{E})$, where $(u, i) \in \mathcal{E}$ denotes an observed interaction and $\mathcal{N}_u, \mathcal{N}_i$ denote the 1-hop graph neighborhoods of user $u$ and item $i$. Instead of initializing the base embedding matrix $\mathbf{E}^{(0)} \in \mathbb{R}^{(|\mathcal{U}| + |\mathcal{I}|) \times d}$ from a random Xavier uniform distribution, we initialize $\mathbf{e}_u^{(0)}$ and $\mathbf{e}_i^{(0)}$ directly with the pre-trained LLM embeddings from Equations~\ref{eq:item_embedding_init} and \ref{eq:user_embedding_init}. As shown in \autoref{fig:recommend_with_lightgcn}, across $L$ graph convolution layers ($l = 0, \dots, L-1$), LightGCN propagates these semantically initialized representations over the bipartite graph using symmetric square-root degree normalization, omitting feature transformation matrices and non-linear activations:}
    {\color{red}\begin{equation}
        \mathbf{e}_u^{(l+1)} = \sum_{i \in \mathcal{N}_u} \frac{1}{\sqrt{|\mathcal{N}_u|}\sqrt{|\mathcal{N}_i|}} \mathbf{e}_i^{(l)}, \qquad
        \mathbf{e}_i^{(l+1)} = \sum_{u \in \mathcal{N}_i} \frac{1}{\sqrt{|\mathcal{N}_i|}\sqrt{|\mathcal{N}_u|}} \mathbf{e}_u^{(l)}.
        \label{eq:lightgcn_propagation}
    \end{equation}}
    \textcolor{red}{After $L$ propagation hops, the layer-wise representations are aggregated via uniform layer combination ($\alpha_l = \frac{1}{L+1}$), and the preference score $\hat{y}_{u,i}^{\text{LightGCN}}$ for user $u$ and candidate item $i$ is computed as the inner product of their final representations:}
    {\color{red}\begin{equation}
        \mathbf{e}_u = \frac{1}{L + 1} \sum_{l=0}^{L} \mathbf{e}_u^{(l)}, \qquad
        \mathbf{e}_i = \frac{1}{L + 1} \sum_{l=0}^{L} \mathbf{e}_i^{(l)}, \qquad
        \hat{y}_{u,i}^{\text{LightGCN}} = \mathbf{e}_u^\top \mathbf{e}_i.
        \label{eq:lightgcn_prediction}
    \end{equation}}
    \textcolor{red}{The embedding parameters $\mathbf{E}^{(0)}$ are fine-tuned end to end using the pairwise Bayesian Personalized Ranking (BPR) loss \cite{bpr} over observed positive pairs $(u, i^+) \in \mathcal{E}$ and sampled unobserved items $(u, i^-)$:}
    {\color{red}\begin{equation}
        \mathcal{L}_{\text{BPR}} = -\sum_{(u, i^+, i^-)} \ln \sigma\left(\hat{y}_{u,i^+}^{\text{LightGCN}} - \hat{y}_{u,i^-}^{\text{LightGCN}}\right) + \lambda \|\mathbf{E}^{(0)}\|_2^2.
        \label{eq:lightgcn_bpr_loss}
    \end{equation}}
\end{itemize}

  \begin{figure}[ht!]
        \centering
        \includegraphics[width=0.9\linewidth]{figures/recommend_with_lightgcn.png}
        \caption{Generating a prediction score for a user’s preference toward an item using a GNN-based architecture.}
        \label{fig:recommend_with_lightgcn}
    \end{figure}

For both the KNN-based and the graph-based approach, we use the embeddings generated by various large language models, and in the experiments we report our findings on which models perform the best.
The hyperparameters used for each model and dataset combination are detailed in the Appendix \textcolor{red}{(\autoref{tab:best_hyperparameters})}.